\documentclass[11pt,a4paper]{article}
\usepackage[margin=0.85in,top=1in,bottom=1in]{geometry}
\usepackage{amsmath,amssymb,amsfonts}
\usepackage{graphicx}
\usepackage{float}
\usepackage{enumitem}
\usepackage{microtype}
\usepackage{caption}
\usepackage[hidelinks]{hyperref}
\usepackage[most,breakable]{tcolorbox}
\tcbuselibrary{breakable,skins}
\usepackage{fontspec}
\setmainfont{DejaVu Serif}
\setsansfont{DejaVu Sans}
\setmonofont{DejaVu Sans Mono}
\captionsetup{font=small,labelfont=bf,skip=4pt}
\setlist{leftmargin=1.5em,topsep=2pt}

%% Breakable problem card — prevents content cutoff across pages
\newtcolorbox{solcard}[3]{
  breakable, enhanced,
  colback=white, colframe=gray!50, arc=2pt,
  title={\textbf{#1}\hfill\textsf{\small #3\,pts}},
  fonttitle=\normalsize\sffamily\bfseries,
  before upper={\small\textit{#2}\par\smallskip\hrule height 0.4pt\smallskip},
  top=4pt, bottom=6pt, left=7pt, right=7pt,
  parbox=false,
}

\title{\textbf{Off-center ball movements --- Solution}}
\author{\normalsize X24 (2024) — English translation}
\date{}
\begin{document}\maketitle\vspace{-0.5em}
\begin{solcard}{A1}{Express the velocity component $\vec{u}_{A}$ of point $A$ in terms of the velocity component $\vec{u}_C$ of the center of the ball, its angular velocity $\vec{\omega}$, as well as the radius vector $\vec{r}$ at an arbitrary moment. Also obtain the time derivative $\dot{\vec{u}}_A$ of the vector $\vec{u}_A$. Express your answer using $\dot{\vec{u}}_C$, $\dot{\vec{\omega}}$ and $\vec{r}$.}{0.40}

The speed of point $A$ is given by the expression: $$\vec{v}_A=\vec{v}_C+\bigl[\vec{\omega}\times\vec{r}\bigr]{.} $$Since vector $\vec{r}$ is perpendicular to the plane of the wall, we have:


Since the wall is flat, the vector $\vec{r}$ remains constant during the movement, therefore, after differentiation we have:

\par\smallskip\noindent\textbf{Answer:} \textbackslash{}textbf{Answer:} $$\vec{u}_A=\vec{u}_C+\bigl[\vec{\omega}\times\vec{r}\bigr]{.} $$
\end{solcard}

\begin{solcard}{A2}{Determine the friction force $\vec{F}_0$ acting on the ball at the initial moment of contact with the wall. Express your answer in terms of $\vec{e}_x$, $\vec{e}_z$, $\alpha$, $\mu$ and the normal reaction force of the wall $N_0$ at the initial moment.}{0.60}

At the initial moment we have: $$\vec{u}_C(0)=v\sin\alpha\vec{e}_x\qquad \bigl[\vec{\omega}(0)\times\vec{r}\bigr]=-v\cos\alpha\vec{e}_z\Rightarrow \vec{u}_A(0)=v\left(\sin\alpha\vec{e}_x-\cos\alpha\vec{e}_z\right){.} $$Note that vector $(\sin\alpha\vec{e}_x-\cos\alpha\vec{e}_z)$ является единичным. Friction force $\vec{F}$ equals по модулю $\mu N$ and направлена противоположно компоненте скорости $\vec{u}_A$, therefore: $$\vec{F}=-\cfrac{\vec{u}_A}{|\vec{u}_A|}\cdot \mu N{,} $$or:

\par\smallskip\noindent\textbf{Answer:} \textbackslash{}textbf{Answer:} $$\vec{F}=\mu N(\cos\alpha\vec{e}_z-\sin\alpha\vec{e}_x){.} $$
\end{solcard}

\begin{solcard}{A3}{Prove that the time derivative $\dot{\vec{u}}_A$ of the velocity component $\vec{u}_A$ is related to the friction force $\vec{F}$ by the relation: $$\dot{\vec{u}}_A=\cfrac{7\vec{F}}{2m}{.} $$This fact can be used further, even if you could not prove it.}{1.00}

From the theorem on the motion of the center of mass for a ball it follows: $$m\dot{\vec{u}}_C=\vec{F}{.} $$Let $I=2mR^2/5$ – moment of inertia однородного ball/sphereа относительно diameterа. Since ball/sphere сферически симметричен, его angular momentum $\vec{L}_C$ относительно центра масс составляет: $$\vec{L}_C=I\vec{\omega}{.} $$Let us write основное equation динамики вращательного движения относительно центра масс ball/sphereа: $$\cfrac{d\vec{L}_C}{dt}=I\dot{\vec{\omega}}=\vec{M}=\bigl[\vec{r}\times\vec{F}\bigr]{.} $$Умножая vectorно from the left на $\vec{r}⟦ 8⟧$I\textbackslash{}bigl[\textbackslash{}vec{r}\textbackslash{}times\textbackslash{}dot{\textbackslash{}vec{\textbackslash{}omega}}\textbackslash{}bigr]=\textbackslash{}bigl[\textbackslash{}vec{r}\textbackslash{}times\textbackslash{}bigl[\textbackslash{}vec{r}\textbackslash{}times\textbackslash{}vec{F}\textbackslash{}bigr]\textbackslash{}bigr]=\textbackslash{}vec{r}\textbackslash{}bigl(\textbackslash{}vec{r}\textbackslash{}cdot\textbackslash{}vec{F}\textbackslash{}bigr)-r^2\textbackslash{}vec{F}=-r^2\textbackslash{}vec{F}\textbackslash{}Rightarrow \textbackslash{}vec{F}=\textbackslash{}cfrac{I\textbackslash{}bigl[\textbackslash{}dot{\textbackslash{}vec{\textbackslash{}omega}}\textbackslash{}times\textbackslash{}vec{r}\textbackslash{}bigr]}{r^2}{.} $$В последнем переходе мы учли, что $\vec{r}\perp\vec{F}$. Воспользуемся результатом пункта $\mathrm{A1}$: $$\textbackslash{}dot{\textbackslash{}vec{u}}_A=\textbackslash{}dot{\textbackslash{}vec{u}}_C+\textbackslash{}bigl[\textbackslash{}dot{\textbackslash{}vec{\textbackslash{}omega}}\textbackslash{}times\textbackslash{}vec{r}\textbackslash{}bigr]=\textbackslash{}cfrac{\textbackslash{}vec{F}}{m}+\textbackslash{}cfrac{\textbackslash{}vec{F}r^2}{I}=\textbackslash{}cfrac{7\textbackslash{}vec{F}}{2m}{.} $$

\end{solcard}

\begin{solcard}{A4}{Determine the velocity component $\vec{u}_{Aк}$ immediately after the collision, assuming that the ball slides along the wall during the entire time of the collision. Express your answer in terms of $v$, $\alpha$, $\mu$, $\vec{e}_x$ and $\vec{e}_z$. At what maximum value of the friction coefficient $\mu_{max}$ does slippage not stop during the entire impact time? Express your answer using $\alpha$.}{0.50}

The friction force vector $\vec{F}$ is directed against the velocity component $\vec{u}_A$ and codirected with the derivative of the velocity component $\dot{\vec{u}}_A$. It follows from this that the directions $\vec{u}_A$ and $\vec{F}$ are preserved during the collision. Then from point $\mathrm{A3}$ we have: $$\cfrac{du_A}{dt}=-\cfrac{7F}{2m}=-\cfrac{7\mu N}{2m}\Rightarrow u_{A\text{к}}-u_A(0)=u_{A\text{к}}-v=-\cfrac{7\mu}{2m}\int\limits_{0}^t Ndt{.} $$До соударения компонента скорости центра ball/sphereа $v_{Cy}(0)$ equals $-v\cos\alpha$. Since удар упругий, ссу после соударения компонента скорости центра ball/sphereа $v_{Cy}$ equals $v\co s\alpha$. Then for momentumа силы реакции $N$ we have: $$\int\limits_{0}^tNdt=\Delta{p}_y=2mv\cos\alpha{.} $$Thus: $$u_{A\text{к}}=v(1-7\mu\cos\alpha){,} $$Проскальзывание не прекращается, if $u_{A\text{к}}\geq 0$, что for/atводит к следующему ограничению на $\mu$: $$\mu\leq\cfrac{1}{7\cos\alpha}{.} $$If however $\mu\geq 1/(7\cos\alpha)$, то $u_{A\text{к}}=0$. Finally:

\par\smallskip\noindent\textbf{Answer:} \textbackslash{}textbf{Answer:} $$\vec{u}_{A\text{к}}= \begin{cases} v(1-7\mu\cos\alpha)(\sin\alpha\vec{e}_x-\cos\alpha\vec{e}_z)\quad\text{при}\quad\mu\leq\cfrac{1}{7\cos\alpha}\\ 0\quad\text{при}\quad\mu\geq\cfrac{1}{7\cos\alpha} \end{cases} $$
\end{solcard}

\begin{solcard}{A5}{At $\mu<\mu_{max}$, determine the speed of the center of the ball $\vec{v}_{Cк}$, as well as at what angle $\beta$ to the horizon it is directed immediately after the collision. Express your answers using $v$, $\alpha$, $\mu$, $\vec{e}_x$, $\vec{e}_y$ and $\vec{e}_z$.}{0.60}

For angle $\beta$ we have: $$\beta=\operatorname{arctg}\cfrac{v_{Cz}}{\sqrt{v^2_{Cx}+v^2_{Cy}}}{.} $$После удара vector скорости центра ball/sphereа составляет: $$\vec{v}_C=v\cos\alpha\vec{e}_y+\vec{u}_C{.} $$For $\vec{u}_C$ we have: $$\vec{u}_C=\vec{u}_C(0)+\int\limits_0^t\cfrac{\vec{F}dt}{m}=v\sin\alpha\vec{e}_x+\cfrac{2(\vec{u}_A-\vec{u}_A(0))}{7} $$For/At $\mu\leq 1/(7\cos\alpha)$ we have: $$\vec{u}_C=v\sin\alpha\vec{e}_x-\cfrac{2}{7}\cdot 7\mu v\cos\alpha(\sin\alpha\vec{e}_x-\cos\alpha\vec{e}_z)=v\sin\alpha(1-2\mu\cos\alpha)\vec{e}_x+2\mu v\cos^2\alpha\vec{e}_z{.} $$Thus: $$\vec{v}_C=v\sin\alpha(1-2\mu\cos\alpha)\vec{e}_x+v\cos\alpha\vec{e}_y+2\mu v\cos^2\alpha\vec{e}_z{,} $$or: $$\beta=\operatorname{arctg}\cfrac{2\mu\cos^2\alpha}{\sqrt{\cos^2\alpha+\sin^2\alpha(1-2\mu\cos\alpha)^2}}{.} $$If however $\mu\geq 1/(7\cos\alpha)$ – $\vec{u}_A=0$. Then: $$\vec{u}_C=v\sin\alpha\vec{e}_x-\cfrac{2v(\sin\alpha\vec{e}_x-\cos\alpha\vec{e}_z)}{7}=\cfrac{5v\sin\alpha\vec{e}_x}{7}+\cfrac{2v\cos\alpha\vec{e}_z}{7}{.} $$Thus: $$\vec{v}_C=\cfrac{5v\sin\alpha\vec{e}_x}{7}+v\cos\alpha\vec{e}_y+\cfrac{2v\cos\alpha\vec{e}_z}{7}{,} $$or: $$\beta=\operatorname{arctg}\cfrac{2\cos\alpha}{\sqrt{(5\sin\alpha)^2+(7\cos\alpha)^2}}{.} $$Final:

\par\smallskip\noindent\textbf{Answer:} \textbackslash{}textbf{Answer:} $$\beta=\begin{cases} \operatorname{arctg}\cfrac{2\mu\cos^2\alpha}{\sqrt{\cos^2\alpha+\sin^2\alpha(1-2\mu\cos\alpha)^2}}\quad\text{при}\quad \mu\leq\cfrac{1}{7\cos\alpha}\\ \operatorname{arctg}\cfrac{2\cos\alpha}{\sqrt{(5\sin\alpha)^2+(7\cos\alpha)^2}}\quad\text{при}\quad \mu\geq\cfrac{1}{7\cos\alpha} \end{cases} $$
\end{solcard}

\begin{solcard}{A6}{At $\mu<\mu_{max}$, determine the coordinates $x_C$, $y_C$ of the center of the ball at the moment it falls on the table. Express your answers using $v$, $g$, $\mu$ and $\alpha$.}{0.40}

The ball will fall on the table surface after time $t$, equal to: $$t=\cfrac{2v_{Cz}}{g}{.} $$Then for координат $x_C$ and $y_C$ we have: $$x_C=v_{Cx}t=\cfrac{2v_{Cx}v_{Cz}}{g}\qquad y_C=v_{Cy}t=\cfrac{2v_{Cy}v_{Cz}}{g}{.} $$Substituting $v_{Cx}$, $v_{Cy}$ and $v_{Cz}$ for/at сных значениях $\mu$, we get:

\par\smallskip\noindent\textbf{Answer:} \textbackslash{}textbf{Answer:} $$x_C=\begin{cases} \cfrac{4\mu v^2\cos^2\alpha\sin\alpha(1-2\mu\cos\alpha)}{g}\quad\text{при}\quad \mu\leq\cfrac{1}{7\cos\alpha}\\ \cfrac{10v^2\sin 2\alpha}{49g}\quad\text{при}\quad \mu\geq\cfrac{1}{7\cos\alpha} \end{cases} $$
\end{solcard}

\begin{solcard}{A7}{For arbitrary values ​​of $\mu$, determine the amount of heat $Q$ released during the collision of the ball with the wall. Express your answer using $m$, $v$, $\mu$ and $\alpha$. Note: explicit calculation of the work done by the friction force will significantly simplify the solution of the problem.}{1.00}

Note: explicit calculation of the work done by the friction force will significantly simplify the solution of the problem.


For the power $P_F$ of the friction force we have: $$P_F=\vec{F}\cdot\vec{u}_A=\cfrac{2m\vec{u}_A\cdot\dot{\vec{u}}_A}{7}{.} $$Then for количества теплоты $Q$ we have: $$Q=-\int\limits_{0}^{t}P_Fdt=\int\limits_{u_{A\text{к}}}^{u_{A}(0)}\cfrac{2mu_A\dot{u}_A}{7}=\cfrac{m(u^2_A(0)-u^2_{A\text{к}})}{7}{.} $$Taking into account that $u_A(0)=v$, а also expression for $u_{A\text{к}}(\mu)$, we obtain:

\par\smallskip\noindent\textbf{Answer:} \textbackslash{}textbf{Answer:} $$Q=\begin{cases} mv^2(2\mu\cos\alpha-7\mu^2\cos^2\alpha)\quad\text{при}\quad \mu\leq\cfrac{1}{7\cos\alpha}\\ \cfrac{mv^2}{7}\quad\text{при}\quad \mu\geq\cfrac{1}{7\cos\alpha} \end{cases} $$
\end{solcard}

\begin{solcard}{B1}{Determine the velocity vector components of the ball's center $v_\varphi$ and $v_z$ in a cylindrical coordinate system. Express your answers using $r$, $\dot\varphi$ and $\dot{z}$.}{0.20}

The radius vector $\vec{r}_C$ of the center of the ball can be represented in the following form: $$\vec{r}_C=r\vec{e}_r+z\vec{e}_z{.} $$Differentiating: $$\vec{v}_C=\cfrac{d\vec{r}_C}{dt}=r\dot{\vec{e}}_r+\dot{z}\vec{e}_z=r\dot{\varphi}\vec{e}_\varphi+\dot{z}\vec{e}_z{.} $$Thus:

\par\smallskip\noindent\textbf{Answer:} \textbackslash{}textbf{Answer:} $$v_\varphi=r\dot{\varphi}\qquad v_z=\dot{z}{.} $$
\end{solcard}

\begin{solcard}{B2}{Determine the acceleration vector components of the ball's center $a_r$, $a_\varphi$ and $a_z$ in a cylindrical coordinate system. Express your answers using $r$, $v_\varphi$, $\dot{v}_\varphi$ and $\dot{v}_z$.}{0.30}

Let us differentiate the velocity vector with respect to time: $$\vec{a}_C=\cfrac{d\vec{v}_C}{dt}=v_\varphi\dot{\vec{e}}_\varphi+\dot{v}_\varphi\vec{e}_\varphi+\dot{v}_z\vec{e}_z=-v_\varphi\dot{\varphi}\vec{e}_r+\dot{v}_\varphi\vec{e}_\varphi+\dot{v}_z\vec{e}_z{.} $$Taking into account that $v_\varphi=r\dot{\varphi}$, we obtain:

\par\smallskip\noindent\textbf{Answer:} \textbackslash{}textbf{Answer:} $$a_r=-\cfrac{v^2_\varphi}{r}\qquad a_\varphi=\dot{v}_\varphi\qquad a_z=\dot{v}_z{.} $$
\end{solcard}

\begin{solcard}{B3}{From the condition of no slipping, determine the components of the angular velocity of the ball $\omega_\varphi$ and $\omega_z$ in the cylindrical coordinate system. Express your answers using $r$, $v_\varphi$ and $v_z$.}{0.40}

For the speed of the point $A$ of the ball at which it contacts the edge of the table, we have: $$\vec{v}_A=\vec{v}_C+\bigl[\vec{\omega}\times\overrightarrow{CA}\bigr]{,} $$where $\overrightarrow{CA}$ – vector, проведённый от центра $C$ ball/sphereа в точку $A$. Let us write condition equalsства нулю скорости точки $A$ ball/sphereа, в которой он контактирует с краем стола: $$v_{A\varphi}=v_\varphi-\omega_zr\qquad v_{Az}=v_z+\omega_\varphi r{.} $$Thus:

\par\smallskip\noindent\textbf{Answer:} \textbackslash{}textbf{Answer:} $$\omega_z=\cfrac{v_\varphi}{r}\qquad \omega_\varphi=-\cfrac{v_z}{r}{.} $$
\end{solcard}

\begin{solcard}{C1}{Determine the friction force $F_\varphi(\varphi)$ component acting on the ball, as well as the acceleration component $a_\varphi(\varphi)$ of its center. Express your answers in terms of the mass of the ball $m$, $g$ and $\varphi$.}{0.80}

Let us write down the theorem on the motion of the center of mass for a ball in projection onto the axis directed along the vector $e_\varphi$: $$ma_\varphi=F_\varphi+mg\sin\varphi{.} $$Write the equation динамики вращательного движения относительно оси $z$, проходящей via/through центр ball/sphereа: $$I_C\varepsilon_z=I_C\dot{\omega}_z=-rF_\varphi{.} $$Since $\dot{\omega}_z=\dot{v}_\varphi/r=a_\varphi/r$, после деления уравнений we get: $$\cfrac{mr^2}{I_C}=-\cfrac{F_\varphi+mg\sin\varphi}{F_\varphi}{,} $$from:


Excluding $F_\varphi$, we get:

\par\smallskip\noindent\textbf{Answer:} \textbackslash{}textbf{Answer:} $$F_\varphi=-\cfrac{mg\sin\varphi}{1+mr^2/I_C}=-\cfrac{2mg\sin\varphi}{7}{.} $$
\end{solcard}

\begin{solcard}{C2}{Get addicted $v_\varphi(\varphi)$. Express your answer in terms of $v$, $g$, $r$, $\alpha$ and $\varphi$.}{0.50}

Multiplying the expression for $v_\varphi$ and taking into account that $v_\varphi=r\dot{\varphi}$, we get: $$a_\varphi v_\varphi=v_\varphi\dot{v}_\varphi=\cfrac{5gv_\varphi\sin\varphi}{7}=\cfrac{5gr\sin\varphi\dot{\varphi}}{7}{.} $$Let us integrate поrayенное expression по времени: $$\int\limits_{v\cos\alpha}^{v_\varphi(\varphi)}v_\varphi dv_\varphi=\int\limits_{0}^{\varphi}\cfrac{5gr\sin\varphi d\varphi}{7}\Rightarrow \cfrac{v^2_\varphi-v^2\cos^2\alpha}{2}=\cfrac{5gr(1-\cos\varphi)}{7}{.} $$Thus:

\par\smallskip\noindent\textbf{Answer:} \textbackslash{}textbf{Answer:} $$v_\varphi=\sqrt{v^2\cos^2\alpha+\cfrac{10gr(1-\cos\varphi)}{7}}{.} $$
\end{solcard}

\begin{solcard}{C3}{Under what condition does the ball not come off the table at the moment when the lowest point of the ball reaches its edge? Write this condition using $v$, $g$, $r$ and $\alpha$. In all subsequent paragraphs, assume that this condition is met.}{0.20}

For the reaction force $N$ at the moment when the lower point of the ball reaches the edge of the table, for the reaction force $N$ we have: $$N=mg-\cfrac{mv^2\cos^2\alpha}{r}{.} $$Отрывnotт for/at условии $N\geq{0}$, therefore:

\par\smallskip\noindent\textbf{Answer:} \textbackslash{}textbf{Answer:} $$v\cos\alpha\leq\sqrt{gr}{.} $$
\end{solcard}

\begin{solcard}{C4}{Determine the angle $\varphi_1$ at the moment the ball leaves the table. Express your answer in terms of $v$, $g$, $r$ and $\alpha$.}{0.50}

At the moment the ball comes off the table $N=0$, therefore we have: $$N=mg\cos\varphi_1-\cfrac{mv^2_{\varphi_1}}{r}{.} $$For/Atравняем два выражения for $v^2_\varphi$: $$v^2_{\varphi_1}=gr\cos\varphi_1=v^2\cos^2\alpha+\cfrac{10gr(1-\cos\varphi_1)}{7}{.} $$Thus: $$\cos\varphi_1=\cfrac{10gr+7v^2\cos^2\alpha}{17gr}{,} $$hence for $\varphi_1$ we find:

\par\smallskip\noindent\textbf{Answer:} \textbackslash{}textbf{Answer:} $$\varphi_1=\arccos\left(\cfrac{10}{17}+\cfrac{7v^2\cos^2\alpha}{17gr}\right){.} $$
\end{solcard}

\begin{solcard}{D1}{Express the kinetic energy of the ball $E_k$ in terms of $m$, $v_\varphi$, $v_z$, $\omega_r$ and $r$.}{0.50}

Let's use Koenig's theorem: $$E_k=\cfrac{mv^2_C}{2}+\cfrac{I_C\omega^2}{2}=\cfrac{mv^2_\varphi}{2}+\cfrac{mv^2_z}{2}+\cfrac{I_C\omega^2_r}{2}+\cfrac{I_C\omega^2_\varphi}{2}+\cfrac{I_C\omega^2_z}{2}{.} $$Substituting $\omega_\varphi$ and $\omega_z$, we get:

\par\smallskip\noindent\textbf{Answer:} \textbackslash{}textbf{Answer:} $$E_k=\left(m+\cfrac{I_C}{r^2}\right)\cfrac{v^2_\varphi}{2}+\left(m+\cfrac{I_C}{r^2}\right)\cfrac{v^2_z}{2}+\cfrac{I_C\omega^2_r}{2}=\cfrac{7m(v^2_\varphi+v^2_z)}{10}+\cfrac{mr^2\omega^2_r}{5}{.} $$
\end{solcard}

\begin{solcard}{D2}{Write down the law of conservation of mechanical energy for the ball. Combining it with the result of item $\mathrm{C2}$, show that the quantities $\omega_r$ and $v_z$ are related by the relation: $$1=\cfrac{\omega^2_r}{A^2}+\cfrac{v^2_z}{B^2}{,} $$where $A,B>0$ — постоянные coefficientы. Determine $A$ and $B$. Answerы выtimesandте via/through $v$, $r$ and $\alpha$.}{0.60}

The law of conservation of energy is as follows: $$E_k=\cfrac{7m(v^2_\varphi+v^2_z)}{10}+\cfrac{mr^2\omega^2_r}{5}=E_{k0}+A_\text{тяж}=\cfrac{7mv^2}{10}+mgr(1-\cos\varphi){.} $$hence: $$v^2_\varphi=v^2+\cfrac{10gr(1-\cos\varphi)}{7}-v^2_z-\cfrac{2\omega^2_rr^2}{7}{.} $$Substituting зависимость $v^2_\varphi(\varphi)$, we get: $$v^2\cos^2\alpha+\cfrac{10gr(1-\cos\varphi)}{7}=v^2+\cfrac{10gr(1-\cos\varphi)}{7}-v^2_z-\cfrac{2\omega^2_rr^2}{7}{,} $$hence: $$v^2_z+\cfrac{2\omega^2_rr^2}{7}=v^2\sin^2\alpha\Rightarrow 1=\cfrac{v^2_z}{v^2\sin^2\alpha}+\cfrac{2\omega^2_rr^2}{7v^2\sin^2\alpha}{.} $$For $A$ and $B$ we have:

\par\smallskip\noindent\textbf{Answer:} \textbackslash{}textbf{Answer:} $$A=\sqrt{\cfrac{7}{2}}\cfrac{v\sin\alpha}{r}\qquad B=v\sin\alpha{.} $$
\end{solcard}

\begin{solcard}{D3}{The angular acceleration vector $\vec{\varepsilon}$ of the ball can be represented as: $$\vec{\varepsilon}=\varepsilon_r\vec{e}_r+\varepsilon_\varphi\vec{e}_\varphi+\varepsilon_z\vec{e}_z{.} $$Using equation динамики вращательного движения относительно центра ball/sphereа, покажите, что $\varepsilon_r=0$. Using поrayенное equality, выtimesandте $\dot{\omega}_r$ via/through $\dot{\varphi}$, $v_z$ and $r$.}{0.50}

From the equation of dynamics of rotational motion relative to the center of the ball we have: $$\cfrac{d\vec{L}_C}{dt}=I\vec{\varepsilon}=\bigl[\overrightarrow{CA}\times\left(\vec{F}+\vec{N}\right)\bigr]{.} $$Since $\vec{N}\parallel\overrightarrow{CA}$, we have: $$I\vec{\varepsilon}=\bigl[\overrightarrow{CA}\times\vec{F}\bigr]{.} $$


Let's use the expression for angular acceleration in a cylindrical coordinate system: $$\varepsilon_r=\dot{\omega}_r-\dot{\varphi}\omega_\varphi=0{,} $$from:

\par\smallskip\noindent\textbf{Answer:} \textbackslash{}textbf{Answer:} Since $\vec{\varepsilon}\perp\overrightarrow{CA}$, the angular acceleration component is $\varepsilon_r=0$.
\end{solcard}

\begin{solcard}{D4}{Combining the results of points $\mathrm{D2}$ and $\mathrm{D3}$, you obtain dependencies $\omega_r(\varphi)$ and $v_z(\varphi)$. Express your answers using $v$, $\alpha$, $r$ and $\varphi$.}{1.20}

Let's use the result of point $\mathrm{B4}$: $$\dot{\omega}_r=-\cfrac{\dot{\varphi}v_z}{r}{.} $$Expressing $v_z$ and substituting в equation, поrayенное в предыдущем partе, we find: $$1=\cfrac{\omega^2_r}{A^2}+\cfrac{r^2\dot{\omega}^2_r}{B^2\dot{\varphi}^2}{.} $$Note that это equation с сделяющимися переменными $\omega_r$ and $\varphi$: $$d\varphi=-\cfrac{r}{B}\cfrac{d\omega_r}{\sqrt{1-\cfrac{\omega^2_r}{A^2}}}{.} $⟦ 9⟧\dot{\omega}_r{<}0$, а значит и $\omega_r{<}0$. Вводя переменную $t=\o mega_r/A$ и интегрируя, находим: $$\textbackslash{}varphi=-\textbackslash{}cfrac{rA}{B}\textbackslash{}int\textbackslash{}limits_{0}^{t(\textbackslash{}varphi)}\textbackslash{}cfrac{dt}{\textbackslash{}sqrt{1-t^2}}=-\textbackslash{}cfrac{rA}{B}\textbackslash{}arcsin t\textbackslash{}Biggl|_{0}^{t(\textbackslash{}varphi)}\textbackslash{}Rightarrow t(\textbackslash{}varphi)=-\textbackslash{}sin\textbackslash{}left(\textbackslash{}cfrac{B\textbackslash{}varphi}{rA}\textbackslash{}right){,} $$или же: $$\textbackslash{}omega_r(\textbackslash{}varphi)=-A\textbackslash{}sin\textbackslash{}left(\textbackslash{}cfrac{B\textbackslash{}varphi}{rA}\textbackslash{}right){.} $$Подставляя $t(\varphi)$ в уравнение, связывающее $v_ z$ и $\omega_r$, находим: $$\textbackslash{}cfrac{v^2_z}{B^2}=1-t^2=\textbackslash{}cos^2\textbackslash{}left(\textbackslash{}cfrac{B\textbackslash{}varphi}{rA}\textbackslash{}right){,} $$или же: $$v_z=B\textbackslash{}cos\textbackslash{}left(\textbackslash{}cfrac{B\textbackslash{}varphi}{rA}\textbackslash{}right){,} $$поскольку $v_z(0){>}0$. Finally:

\par\smallskip\noindent\textbf{Answer:} \textbackslash{}textbf{Answer:} $$\omega_r(\varphi)=-\sqrt{\cfrac{7}{2}}\cfrac{v\sin\alpha}{r}\sin\left(\sqrt{\cfrac{2}{7}}\varphi\right)\qquad v_z=v\sin\alpha\cos\left(\sqrt{\cfrac{2}{7}}\varphi\right){.} $$
\end{solcard}

\begin{solcard}{D5}{Let us consider the passage to the limit when the angle is $\alpha\to\pi/2$, i.e. the movement of the ball until it contacts the edge of the table occurs almost parallel to it. Determine the projection of the velocity $v_z$ of the center of the ball, as well as the projection of its angular velocity $\omega_y$ onto the axis $y$, directed vertically downward, at the moment the ball lifts off the table. Express your answers using $v$ and $r$. All numerical coefficients in the answer must be analytical, not approximate!}{0.80}

At $\alpha=\pi/2$ for $\varphi_1$ we have: $$\varphi_1=\arccos\left(\cfrac{10}{17}\right){.} $$


The projection of the angular velocity of the ball onto the axis $y$, directed vertically downward, is equal to: $$\omega_y=\omega_\varphi\sin\varphi_1-\omega_r\cos\varphi_1{.} $$Substituting $\omega_r$, $\omega_\varphi$, we get: $$\omega_y=\cfrac{v}{r}\left(\sqrt{\cfrac{7}{2}}\cos\varphi_1\sin\left(\sqrt{\cfrac{2}{7}}\varphi_1\right)-\sin\varphi_1\cos\left(\sqrt{\cfrac{2}{7}}\varphi_1\right)\right){.} $$После подстановки $\varphi_1$ we find:

\par\smallskip\noindent\textbf{Answer:} \textbackslash{}textbf{Answer:} Then for speed $v_z$ we find: $$v_z=v\cos\left(\sqrt{\cfrac{2}{7}}\arccos\left(\cfrac{10}{17}\right)\right){.} $$
\end{solcard}

\end{document}